% literature-review.tex — \input dari report.tex (berbagi bibliografi)
\chapter{Tinjauan Pustaka}
\label{ch:litreview}

Bab ini meninjau karya terdahulu di seluruh lapis yang dijembatani dan menegakkan celah
penelitian. Bab ini merupakan hasil konsolidasi penelusuran sistematis (lima sudut,
24 sumber, 118 klaim diekstrak, 25 diverifikasi secara adversarial melalui voting
mayoritas; 15 terkonfirmasi, 10 gugur). Vonis kebaruan berkeyakinan \emph{sedang}---dan
\emph{tinggi} untuk celah inti (koordinasi) setelah verifikasi Fase~0: ketiadaan
dalam pustaka tidak dapat dibuktikan, indeks penelusuran tunggal, satu venue penting
(ACM~DL) berbayar, dan sepuluh verifikasi \emph{abstain} (Bagian~\ref{sec:lit-batas}). Karya
terdekat sebaiknya diperiksa ulang di Google Scholar, DBLP, dan IEEE~Xplore sebelum
pengiriman.

\section{Sinyal Kesehatan dan Semantik \emph{Probe}}
\label{sec:lit-health}
Manajemen kesehatan Kubernetes bersifat \emph{poll}: orkestrator memeriksa kontainer
secara berkala. Roberts dkk.~\cite{roberts2025}, dari komunitas keandalan BEAM,
mengarakterisasi ini sebagai \emph{poll-based container monitoring} (PCM) yang ``memerlukan
penyetelan cermat, dapat menurunkan ketersediaan, dan lambat mendeteksi perubahan
kesehatan,'' lalu mengusulkan \emph{signal-based container monitoring} (SCM), yakni
kontainer memberi sinyal proaktif ke orkestrator. SCM adalah karya terdekat untuk
menurunkan keadaan \emph{probe} dari internal runtime (M1) dan menggantikan \emph{probe}
TCP/HTTP dengan sinyal yang lebih kaya (M9); karena penulisnya berada di ranah Erlang/BEAM,
karya ini wajib dikutip dan dibedakan secara langsung. Kontribusi yang masih terbuka dan
khas BEAM: penurunan semantik \emph{probe} \emph{secara otomatis dari supervision tree} dan
sinyal \emph{komposit} runtime.

\section{Pemulihan Runtime: Terkoordinasi vs.\ Independen}
\label{sec:lit-recovery}
Pustaka keandalan klasik membedakan pemulihan terkoordinasi dan independen serta efek
domino yang menyertainya~\cite{elnozahy2002}. Hal ini melandaskan persoalan ``pemulihan
ganda'' (M2) pada teori mapan, sekaligus menegaskan bahwa unsur barunya adalah model
\emph{delegasi kepemilikan lapis}, bukan pengamatan (klasik) bahwa pemulihan tak
terkoordinasi itu mubazir.

\section{Klastering BEAM di atas Kubernetes}
\label{sec:lit-beam}
Tumpukan dominan adalah libcluster~\cite{libcluster} untuk pembentukan \emph{cluster} dan
Horde~\cite{horde} untuk \emph{supervisor}/\emph{registry} terdistribusi dengan
\emph{handoff}. Horde berbasis CRDT sehingga \emph{eventually consistent}: ia menjamin
ketersediaan dan toleransi partisi tetapi ``tidak dapat menjamin konsistensi.'' Secara
baku, kedua sisi partisi tetap beroperasi; satu-satunya mitigasi bawaan adalah opsi kuorum
yang kasar. \emph{Handoff} state pun tidak otomatis---aplikasi harus menyimpan state di
\texttt{terminate/2} dan memulihkannya via \texttt{handle\_continue/2}, dengan
\emph{shutdown timeout} baku nol---dan, yang krusial, tidak ada kopling ke siklus hidup
\emph{pod} Kubernetes. Bonny~\cite{bonny} menyediakan perkakas matang untuk menulis
\emph{operator} Kubernetes dalam Elixir, sehingga kapabilitas operator di balik M4 sudah
ada di BEAM; yang belum ada adalah jembatan dari keanggotaan/\emph{handoff} ke
\emph{readiness} atau terminasi \emph{pod}. Ketiadaan-ketiadaan inilah bukti utama bahwa
M3, M4, dan M7 baru sebagai kontribusi koordinasi.

\section{\emph{Split-Brain}, \emph{Fencing}, dan \emph{Control Plane} sebagai Arbiter}
\label{sec:lit-splitbrain}
Di JVM, \emph{Split Brain Resolver} Akka menawarkan strategi \texttt{lease-majority} yang
didukung \emph{Kubernetes Lease}~\cite{akkaSBR,akkaLease}: hanya partisi yang berhasil
mengambil \emph{lease} yang boleh bertahan. \emph{Lease} ini memanfaatkan kendali
konkurensi dan konsistensi API \emph{server}/etcd, sehingga \emph{control plane} menjadi
otoritas \emph{fencing} eksternal. Ini adalah karya terdahulu langsung bagi pemakaian
etcd/\emph{Lease} sebagai arbiter (M5) pada \emph{actor runtime}---sehingga \textbf{M5
bukan hal baru sebagai gagasan umum, melainkan hanya sebagai penerapan khas BEAM}, dan
sebaiknya dilebur ke kisah generalisasi. \emph{Control plane} sendiri adalah sistem
konsensus tereplikasi dengan batasan Byzantine yang diketahui~\cite{byzK8s}. (Pada BEAM,
perilaku bawaan terkait adalah resolusi konflik nama gaya \emph{last-writer-wins} pada
\texttt{:global} yang hendak digantikan M5; sejauh mana Horde kehilangan state saat pulih
tidak terkonfirmasi---lihat Bagian~\ref{sec:lit-batas}.)

\section{Pendekatan Lintas-Lapis, Otonomik, dan Formal}
\label{sec:lit-crosslayer}
\emph{Self-healing} orkestrasi sering dibingkai sebagai \emph{control loop} otonomik
MAPE-K~\cite{kephart2003}. Model penyebaran \emph{cloud-native} formal menalar penanganan
kesalahan hanya pada tingkat infrastruktur: model UCC'23~\cite{ucc23} menangani kesalahan
lewat replikasi statis dan pengiriman ulang, serta secara eksplisit menempatkan ``cara
runtime memulihkan diri'' di luar cakupan---menegaskan bahwa model koordinasi
runtime-vs-orkestrator (M2, M6) dan kompilasi sumber tunggal (M6) belum terisi. Sebagai
kontras, Nefele~\cite{nefele} mengusulkan orkestrator proses cloud yang ter-inspirasi
\emph{supervision tree} OTP namun \emph{menggantikan} kontainer/Kubernetes alih-alih
berkoordinasi dengannya---justru menegaskan bahwa pendekatan \emph{koordinasi} (bukan
penggantian) inilah yang masih kosong.

\section{\emph{Actor Runtime} di Kubernetes dan Asimetri Lintas-Platform}
\label{sec:lit-actors}
Akka Management menyediakan pemeriksaan \emph{readiness} berbasis keanggotaan
\emph{cluster} (\texttt{ClusterMembershipCheck}) yang menandai \emph{pod} siap hanya ketika
anggotanya mencapai status \texttt{Up}~\cite{akkaHealth}, dengan \texttt{/ready} dan
\texttt{/alive} yang memang dirancang untuk memetakan ke \emph{probe} Kubernetes. Microsoft
Orleans menyebar di Kubernetes dengan integrasi keanggotaan serupa~\cite{orleans}. Jadi
jembatan \emph{keanggotaan$\rightarrow$readiness} (M4) sudah ada di sisi JVM/.NET---meski
ditulis manual dan berbasis keanggotaan, bukan diturunkan dari \emph{supervision tree}.
Asimetri inilah yang menjadi cerita: tumpukan JVM/Akka punya \emph{plumbing} lintas-lapis,
tumpukan BEAM tidak, dan \emph{tak ada} platform yang memformalkan koordinasi sebagai dua
\emph{control loop} dengan kontrak kepemilikan/\emph{handoff}.

\section{Klaim yang Tidak Lolos Verifikasi}
\label{sec:lit-batas}
Demi kejujuran dan untuk membatasi klaim kebaruan, berikut \emph{tidak} terkonfirmasi:
(i)~bahwa Horde membuang state duplikat saat partisi pulih (hanya ``kedua sisi berlanjut''
yang terkonfirmasi); (ii)~bahwa Bonny \emph{tidak} menyediakan koordinasi lintas-lapis
(tak dapat ditegaskan---hanya bahwa ia menyediakan perkakas operator umum); dan (iii)~bahwa
demo kanonik \texttt{ex\_cluster} mengabaikan \emph{split-brain}/\emph{probe}/\emph{grace}.
Ketiga petunjuk yang sebelumnya \emph{abstain} kini telah diverifikasi manual (Fase~0):
\texttt{akkadotnet-healthcheck} (.NET) \emph{terkonfirmasi} memantulkan keanggotaan
\emph{cluster} ke \emph{probe} Kubernetes---sehingga jembatan keanggotaan$\rightarrow$\emph{readiness}
(M4) ada di JVM \emph{dan} .NET; Beamlens \emph{tidak} memiliki integrasi Kubernetes (alat
diagnostik berbasis LLM di \emph{supervision tree}); dan Akka Coordination-Lease
\emph{terkonfirmasi} memakai CRD (\texttt{akka.io/v1}, \emph{kind} \texttt{Lease}) di etcd
dengan kendali konkurensi optimistik. Penelusuran Scholar/DBLP/IEEE~Xplore tidak menemukan
protokol koordinasi OTP$\leftrightarrow$Kubernetes yang terbit, sehingga celah inti
terkonfirmasi (keyakinan naik dari sedang ke tinggi).

\section{Sintesis: Celah dan Vonis Kebaruan}
\label{sec:lit-synth}
Setiap \emph{primitif} yang mendasari mekanisme kandidat sudah ada---\emph{probe}, PDB, dan
\emph{grace period} Kubernetes; klastering dan \emph{handoff} Horde; \emph{supervision
tree} OTP; kuorum \emph{split-brain}; \emph{fencing} via \emph{Kubernetes Lease}; bahkan
pemantauan kesehatan berbasis sinyal. Yang tidak disediakan sumber mana pun---akademik
maupun perkakas---adalah koordinasi lapis pemulihan runtime (OTP) dan orkestrator
(Kubernetes) sebagai dua \emph{control loop} yang saling berinteraksi dengan protokol
kepemilikan/\emph{handoff} yang jelas. Itulah ruang putih yang dapat dipertahankan.
Tabel~\ref{tab:verdicts} memuat vonis tiap mekanisme (lihat Bab~\ref{ch:mechanisms} untuk
rincian rancangannya); kontribusi baru yang paling bersih adalah \textbf{M3, M7, dan M8}.

\begin{table}[H]
  \centering\small
  \begin{tabular}{@{}cp{0.44\linewidth}p{0.32\linewidth}@{}}
    \toprule
    \# & Mekanisme & Vonis kebaruan \\
    \midrule
    M1 & Semantik \emph{probe} dari \emph{supervision tree} & Sebagian tercakup~\cite{roberts2025} \\
    M2 & Delegasi kepemilikan kegagalan & Baru (model)~\cite{elnozahy2002,ucc23} \\
    M3 & Pengendali \emph{grace}$\leftrightarrow$konvergensi$\leftrightarrow$\emph{probe} & \textbf{Baru} \\
    M4 & \emph{Operator} sadar-BEAM + jembatan kesehatan & Terbelah; jembatan ada di Akka~\cite{akkaHealth} \\
    M5 & Arbiter \emph{split-brain} via Lease/etcd & Sudah ada di Akka~\cite{akkaLease} (terlemah) \\
    M6 & Satu spesifikasi $\rightarrow$ OTP + manifes K8s & Baru~\cite{ucc23} \\
    M7 & PDB/maxUnavailable $\leftrightarrow$ kuorum BEAM & \textbf{Baru}~\cite{k8sPDB} \\
    M8 & \emph{Restart-budget} adaptif $\leftrightarrow$ \emph{backoff} K8s & \textbf{Baru}~\cite{kephart2003} \\
    M9 & Sinyal \emph{liveness} komposit BEAM & Sebagian tercakup~\cite{roberts2025} \\
    \bottomrule
  \end{tabular}
  \caption{Mekanisme kandidat dan vonis kebaruan (keyakinan sedang).}
  \label{tab:verdicts}
\end{table}
